\section{Compositional models and model-based problem solving: proofs}
\label{sec:annex-p16-qualitative-compositional-model-formulation}

The statements below use the reader's explicitly specified relations,
finite spaces and semantics. They do not substitute for the full original
implementation proofs of the systems surveyed by Struss.

\begin{annexitem}{annex:p16.cm-composition}{thm:p16-cm-composition}{Composition, both projection directions and a lost join}
\begin{proof}
An actual global behavior satisfies each local relation on its variables
and the connections by hypothesis, hence satisfies their join. Its actual
observable restriction therefore has that behavior as a consistency witness.
For any set \(A\) of allowed observable tuples,
\(R_a\cap\pi_V^{-1}(A)\ne\varnothing\) holds exactly when
\(\pi_V(R_a)\cap A\ne\varnothing\): a witness projects in one direction,
and membership in the projection supplies a full witness in the other.
Likewise \(R_a\subseteq\pi_V^{-1}(A)\) holds exactly when
\(\pi_V(R_a)\subseteq A\), proving universal-prediction preservation.
The equivalences include empty relations. For the binary example, a joined
tuple has \(z=x=y\); both equal observable pairs extend by taking \(z=x\).
Projecting the two fragments separately permits every \(x\) and every
\(y\), including the unequal pair. This proves the claimed failure.
\end{proof}
\end{annexitem}

\begin{annexitem}{annex:p16.cm-search}{thm:p16-cm-search}{Every candidate, cache pruning and preference order}
\begin{proof}
There are exactly \(N\) assignments in the finite Cartesian product,
including the one empty assignment if there are no components. Enumeration
without repetition and termination of every oracle call give termination
and the stated call bound. Exactness of the oracle makes recorded candidates
precisely the consistent ones. If \(p\) is inconsistent, any complete
extension contains every assertion of \(p\); no such extension can have
a model. Cached-conflict skipping therefore removes no diagnosis. Conversely
an inconsistent candidate is either safely skipped or rejected by the oracle.
Thus every diagnosis and no other candidate is returned.

On the finite returned set, testing whether another member is strictly
preferred selects exactly its maximal elements by definition. In the
order-sensitive version, a maximal diagnosis cannot be skipped by a recorded
dominator. If a nonmaximal diagnosis is considered, a maximal diagnosis above
it exists by finiteness of the preference order and precedes it in the linear
extension. That maximal diagnosis cannot itself have been dominance-skipped
and is already recorded. The nonmaximal candidate can be skipped or removed
by the final maximality test. Without the ordering condition, stopping at the
first survivor or dropping unexamined improvements is not this construction.
An incomplete consistency checker can miss an inconsistent candidate; the
proof then has no exactness premise for accepting that survivor.
\end{proof}
\end{annexitem}

\begin{annexitem}{annex:p16.cm-mode-conflicts}{thm:p16-cm-mode-conflicts}{Complete assignments, kernels and monotone fault families}
\begin{proof}
A diagnosis cannot extend an inconsistent partial assignment, since its
witness would also satisfy that partial theory. Hence it satisfies every
conflict clause. Conversely, if a complete assignment is inconsistent,
its finite set of mode assertions has an inclusion-minimal inconsistent
subset. It falsifies that subset's conflict clause. This proves both
directions, including inconsistent background, whose empty nogood negates
to false and permits no complete diagnosis. A partial assignment has only
diagnostic complete extensions precisely when every such extension satisfies
all these clauses. Taking inclusion-minimal partial terms gives exactly
the kernel definition and its claimed domain-relative implication.

Now use two modes. If every minimal inconsistent partial assignment consists
of normality assertions, a consistent faulty set hits each such support:
otherwise its complete assignment would extend the nogood. Adding faults
continues to hit every support, hence remains consistent by the characterization.
Conversely assume feasible faulty sets are upward closed. If a minimal
nogood contains a fault assertion for \(c\), delete that assertion to
obtain \(p'\). Minimality gives a model of the background with \(p'\).
The exhaustive mode axioms extract a complete diagnostic assignment from
that model. It must have \(c\) normal, since otherwise it extends the
original nogood. Changing just \(c\) to faulty enlarges its faulty set,
so upward closure makes the new assignment diagnostic. But that assignment
extends the original nogood, a contradiction. No minimal nogood contains
a fault assertion. A redundant nogood can add other assertions to an
already inconsistent normality set; the conclusion is consequently about
minimal nogoods, not their arbitrary nonminimal extensions.
\end{proof}
\end{annexitem}

\begin{annexitem}{annex:p16.cm-equal-chain}{thm:p16-cm-equal-chain}{The seventeen-component source example at both abstraction levels}
\begin{proof}
Connections identify each output with the next input. Telescoping all
increments gives final output \(p-q\), where \(p\) counts positive
and \(q\) negative faulty modes. The observed one therefore requires
\(p-q=1\). With \(k=p+q\), this is equivalent to
\(p=(k+1)/2,q=(k-1)/2\). These are nonnegative integers exactly when
\(k\) is odd. Conversely selecting those counts on any odd faulty set and
zero increments elsewhere constructs all intermediate values by cumulative
sums and realizes the observation. Each singleton uses its positive mode,
giving seventeen singleton diagnoses. Toggling one fixed component is a
bijection between even and odd subsets of the seventeen components, so
exactly half of the \(2^{17}\) subsets, namely \(2^{16}\), are feasible.

If a partial assignment has a consistent completion but leaves one component
unassigned, change only that component's increment in the completion to a
different one of \(-1,0,1\). The resulting complete assignment is a valid
mode assignment but its sum is no longer one, so it is not a diagnosis.
The partial assignment thus cannot have every completion diagnostic.
Every kernel specifies all components; a consistent complete assignment
is a kernel because its only completion is itself, and removing any mode
has the just-demonstrated bad completion. For a faulty set of size \(k\)
there are \(\binom{k}{(k+1)/2}\) complete diagnoses, not just one kernel.
This supplies the exact distinction between mode assignments and faulty sets.
\end{proof}
\end{annexitem}

\begin{annexitem}{annex:p16.cm-circuit}{thm:p16-cm-circuit}{The equality graph, all realizations and twenty-one alternatives}
\begin{proof}
Fix an arbitrary set of components required normal. Normal wire equations
give equality edges on the two rails. The observed inactive starter and
dark rear light give cross-rail equality edges when those components are
normal. A normal battery and the observed lit normal head light instead
give inequality edges. Contract all equality edges. An inequality which
becomes a loop is inconsistent. If neither does, the remaining graph has
at most two nonloop inequality edges and is bipartite. Assign zero and one
to its two color classes in each connected component and arbitrary zero
to isolated vertices. Lift those values to the original nodes. This satisfies
every required wire/component equation and the observations, with unrestricted
faulty outputs chosen to match their observed activity. Thus inconsistency
occurs exactly when a battery or head-light inequality contracts to a loop.

Connecting the battery's two rails through equality edges requires both
\(w_1,w_2\) and either the starter's equality or both \(w_3,w_4\)
and the rear light's equality. These are precisely \(K_1,K_2\).
Connecting the head-light rails requires both \(w_5,w_6\) and either
the rear-light equality or both \(w_3,w_4\) and the starter equality,
precisely \(K_3,K_4\). The path characterization follows because the
only cross-rail equality edges are at \(T,R\), and a shortest equality
path need not revisit a vertex. The four supports are incomparable;
removing any member breaks its corresponding equality/inequality obstruction
without containing another support. They are exactly the minimal conflicts.

For minimal hitting sets, if \(R,T\) are both selected, their pair
already hits all four and no other member can appear. Otherwise split
according to selection from the disjoint groups \(A_0,C_0\).
Selecting a member of each gives the nine pairs, and any larger selection
has such a pair as a subset and is nonminimal. Selecting from \(A_0\)
but not \(C_0\) forces \(R\) by \(K_3\); to hit \(K_4\) one
must also select \(T\) or a member of \(D_0\). The first alternative
contains \(\{R,T\}\); minimality in the other leaves precisely
\(\{a,R,d\}\). The symmetric case yields \(\{c,T,d\}\).
Selecting from neither group forces \(T\) by \(K_1\) and \(R\)
by \(K_3\), giving the first pair. Each displayed set hits every
support. Removing a pair member misses respectively an \(A_0\)-based
or \(C_0\)-based support; removing \(a,R,d\)'s members misses
\(K_1,K_3,K_4\), respectively, and the symmetric triple has the
symmetric witnesses. All are minimal. This exhausts the cases and gives
the count \(1+3\cdot3+3\cdot2+3\cdot2=22\).
By the normality-only consistency/hitting characterization, these and
only these are minimal fault localizations of the explicit model.
\end{proof}
\end{annexitem}

\begin{annexitem}{annex:p16.cm-default-sound}{thm:p16-cm-default-sound}{Least default generation prevents a dominated extension}
\begin{proof}
An inconsistent extension would be the entire language. It blocks every
justification, so its generated closure would be \(\operatorname{Th}(W)\),
which is consistent, contradicting the fixed-point condition. Hence the
given extension \(E=\operatorname{Th}(W\cup\{a\})\) is consistent and
\(a\) is a diagnosis.

For a fixed candidate \(E\), extension generation begins with
\(\operatorname{Th}(W)\). A default adds its mode only when its
prerequisite is already entailed and its justification is not contradicted
by \(E\). There are finitely many possible mode conclusions, so this least
generation reaches its fixed point after finitely many additions. Every
added mode belongs to \(E\) and hence agrees with \(a\) by exclusivity.

Suppose a different diagnosis \(a'\) strictly dominates \(a\).
Choose a model \(M'\) of \(W\cup\{a'\}\). If it satisfies every
generated mode conclusion, it satisfies the generated closure \(E\),
and hence \(a\), contradicting exclusivity and \(a'\ne a\).
There must be a first generated mode \(b=a(c)\) false in \(M'\).
All earlier generated conclusions and \(W\) hold in \(M'\).
Its actual mode \(a'(c)\) is different from \(b\) and, by dominance,
strictly preferred to it. The default generating \(b\) required
\(\neg a'(c)\) among its prerequisite conjuncts. That prerequisite was
entailed by \(W\) and the earlier conclusions, all satisfied in \(M'\).
This contradicts \(M'\models a'(c)\). No such dominating diagnosis
exists. The argument proves the stated implication, not its converse.
\end{proof}
\end{annexitem}

\begin{annexitem}{annex:p16.cm-temporal}{thm:p16-cm-temporal}{Restrictions of a trajectory and incompatible snapshot witnesses}
\begin{proof}
A model of a temporal diagnosis supplies a trajectory/history satisfying
the model at every observed time. Restricting that same witness to each
time preserves its state constraints and observed values, proving snapshot
necessity. Conversely the two observations in the example have separate
snapshot witnesses \((x,x')=(0,0)\) and \((1,0)\). All models use the
same derivative/integration semantics and the single persistent mode
\(x'=0\). If a differentiable trajectory joined the observations, the
mean value theorem on \([0,1]\) would give some intermediate derivative
\((x(1)-x(0))/(1-0)=1\), contradicting zero derivative everywhere.
Thus common sequence constraints and a persistent mode do not by themselves
turn independent snapshot witnesses into one full-history witness.
\end{proof}
\end{annexitem}

\begin{annexitem}{annex:p16.cm-trace}{thm:p16-cm-trace}{Every hidden segment and every observed prefix}
\begin{proof}
There are \(|Q|2^{|F|}\) augmented pairs. Hidden closure begins with
the seed set and repeatedly follows a hidden edge, unioning fault labels.
Each insertion has a finite hidden-path witness, inductively concatenating
an edge to its source witness. Conversely induction on the length of any
finite hidden path shows that its target pair is eventually inserted;
once no new pair is found every such edge is closed. Fresh-pair insertion
can occur only the stated finite number of times, even if hidden cycles
can be traversed arbitrarily often.

For an empty observed word, exactly the finite hidden paths from the initial
pair are covered. Suppose the statement holds for a prefix \(w\).
Taking an edge labeled by the next symbol \(z\) after any covered hidden
prefix, then closing under hidden edges, produces only paths whose
projection is \(wz\), with exactly the union of their fault labels.
Conversely any finite path with projection \(wz\) splits at its last
observable edge into a path with projection \(w\), a \(z\)-edge,
and a hidden suffix. Prefix induction, the edge update and hidden-closure
completeness therefore insert its pair. This proves both directions for
all finite words. If some covered pair has empty fault set, its witness
path is a no-fault explanation. If all have nonempty fault sets and there
is at least one pair, every explanation has a fault. With no pair there
is no explanation under the stated model, so fault detection by universal
quantification alone would be vacuous. Infinite unobservable executions
and exact ordered fault histories are outside this statement.
\end{proof}
\end{annexitem}

\begin{annexitem}{annex:p16.cm-discrimination}{thm:p16-cm-discrimination}{The quantifiers, projections and all response-set cases}
\begin{proof}
Applicability gives a nonempty response set for each hypothesis at each
admissible stimulus. An observation can fit both exactly when
\(z\in P_i(u)\cap P_j(u)\); such a response exists exactly when
\(u\in\pi_U(P_i\cap P_j)\). Complementing inside \(A\) therefore
gives all and only the definite stimuli. A response fitting hypothesis
\(i\) and refuting \(j\) exists exactly when
\(u\in\pi_U(P_i\setminus P_j)\). Requiring this for each of the
two hypotheses is intersection of the two projected differences, proving
the formula for \(B_{ij}\). Requiring it for at least one is their union,
proving the formula for \(E_{ij}\). No informative response means neither
difference contains a response, equivalently the nonempty response sets
are identical. Subtracting the definite case or the two-sided case gives
the remaining displayed classes. With nonempty disjoint response sets,
both differences are nonempty, so definite stimuli are indeed in \(B_{ij}\).

For \(\{a\},\{a,b\}\) the shared observation is \(a\), the first
difference is empty and the second is \(\{b\}\): it lies in
\(E_{ij}\), not \(B_{ij}\) or \(D_{ij}\). Thus the complement of
the no-information and definite classes includes it, although the
each-hypothesis quantifier fails. For \(\{a,c\},\{b,c\}\) the
intersection \(\{c\}\) and both differences are nonempty; for
\(\{a\},\{b\}\) the intersection is empty; for identical nonempty
sets both differences are empty. These verify every example and explain
why the source complement cannot establish the stronger characterization.
\end{proof}
\end{annexitem}

\begin{annexitem}{annex:p16.cm-test-sets}{thm:p16-cm-test-sets}{The pair family and minimality in both directions}
\begin{proof}
Pairwise guaranteed discrimination says that for every \(i<j\) some
\(u\in T\) definitely distinguishes that pair. By the preceding
equivalence this is exactly \(T\cap D_{ij}\ne\varnothing\) for
every pair, the hitting-set condition. The feasible families on both sides
are identical, so their inclusion-minimal elements are identical as well.
An empty pair set cannot be hit. With no pairs, every \(T\) meets the
vacuous requirement and only the empty one is minimal. If all pair sets
are nonempty, choosing one stimulus from each produces a finite covering
set, with at most one choice per pair; it need not minimize cost or size.
For any finite feasible set, repeatedly remove a dispensable member until
none remains. Each deletion preserves feasibility and, with nonnegative
additive costs, cannot increase cost. Starting from a cheapest finite
solution therefore gives an inclusion-minimal solution of no greater cost.
Zero costs allow nonminimal solutions to tie; no uniqueness is asserted.
\end{proof}
\end{annexitem}

\begin{annexitem}{annex:p16.cm-information}{thm:p16-cm-information}{Bayes, entropy expansion and the special unknown group}
\begin{proof}
For fixed \(u\), write \(q_i(z)\), \(r(z)\), and
\(p_i(z)=p_iq_i(z)/r(z)\). Work only with terms of positive joint
probability; all others contribute zero. Expanding the posterior entropy
and using \(\log p_i(z)=\log p_i+\log q_i(z)-\log r(z)\) gives
\[
\begin{split}
\sum_z r(z)H(p(\cdot\mid z,u))
 &=-\sum_{i,z}p_iq_i(z)\log p_i
   -\sum_{i,z}p_iq_i(z)\log q_i(z)
   +\sum_{i,z}p_iq_i(z)\log r(z)\\
 &=H(p)+\sum_i p_iH(q_i)-H(r).
\end{split}
\]
The first equality uses the stipulated Bayes probabilities; the second
uses \(\sum_z q_i(z)=1\) and \(\sum_i p_iq_i(z)=r(z)\).
Rearranging proves the information identity. A unique predicted value has
conditional entropy zero, while a uniform distribution on \(m\) values
has entropy \(-m(1/m)\log(1/m)=\log m\). The weighted conditional
entropy is therefore \(p_{unk}\log m\). Negating the information
criterion gives exactly the expression to be minimized.

For the two examples, equal binary outcome probabilities have entropy
\(-2(1/2)\log_2(1/2)=1\). Deterministic hypothesis-specific responses
have conditional entropy zero. Identical fair responses have conditional
entropy one for both hypotheses, and Bayes leaves each posterior at
\((1/2)(1/2)/(1/2)=1/2\). Hence the information is respectively
one and zero bit. No probability for an unknown model is derived from
logical consistency in this calculation.
\end{proof}
\end{annexitem}

\begin{annexitem}{annex:p16.cm-remedies}{thm:p16-cm-remedies}{Actual witnesses, all remaining cases, cost and reachability}
\begin{proof}
A consistent theory including a complete target assignment has a model
of the theory with that assignment removed. Conversely a model of
\(L\cup S\cup\{a\}\cup O_P\cup G\) supplies one state for each
reconfigurable component by the exhaustive exclusive domain axioms.
Collecting these states gives a complete target assignment true in the
same model. This proves the existence equivalence. To see why \(O_P\)
cannot be dropped, take a binary persistent variable \(x\) observed zero
and a goal \(x=1\), with no other restriction. The theory with the
goal alone is consistent, but retaining the actual persistent observation
makes every target theory inconsistent.

For each remaining hypothesis, intersection with \(G_i\) nonempty is
exactly the existence of a permitted goal-satisfying outcome, whereas
nonempty \(R_i(d)\subseteq G_i\) says every permitted outcome satisfies
the goal and the candidate has an outcome. Universal quantification over
the hypotheses yields the displayed two robust criteria. The subset
criterion implies nonempty intersection. For the converse counterexample,
let one candidate permit two voltage outcomes \(\{0,1\}\), with goal
\(\{1\}\). The intersection is nonempty but outcome zero violates
the goal. The example also shows why a weak behavior relation need not
give a physical supply guarantee.

For a finite cheapest feasible intervention/change set \(d\), consider
all its feasible subsets, a nonempty finite family containing \(d\).
Choose an inclusion-minimal one \(d'\). Inclusion-monotone cost gives
\(cost(d')\le cost(d)\), proving the stated result without assuming
feasibility itself is monotone. If the candidate space restricts admissible
sets, minimize within that space. Finally consider a transition graph
with a start vertex and a different goal vertex, each satisfying its
state constraints, but with no edge/path between them. The goal is a
valid state yet unreachable. Intermediate safety restrictions can likewise
remove every path to an otherwise feasible goal; a target-state witness
alone does not prove a plan exists.
\end{proof}
\end{annexitem}
